\begin{lemma}
\label{lemma-silly-normal}
Let $R$ be a ring. Let $x, y \in R$ be nonzerodivisors.
Let $R[x/y] \subset R_{xy}$ be the $R$-subalgebra generated
by $x/y$, and similarly for the subalgebras $R[y/x]$ and $R[x/y, y/x]$.
If $R$ is integrally closed in $R_x$ or $R_y$, then the sequence
$$
0 \to R \xrightarrow{(-1, 1)} R[x/y] \oplus R[y/x] \xrightarrow{(1, 1)}
R[x/y, y/x] \to 0
$$
is a short exact sequence of $R$-modules.
\end{lemma}

\begin{proof}
Since $x/y \cdot y/x = 1$ it is clear that the map
$R[x/y] \oplus R[y/x] \to R[x/y, y/x]$ is surjective.
Let $\alpha \in R[x/y] \cap R[y/x]$. To show exactness in the middle
we have to prove that $\alpha \in R$. By assumption we may write
$$
\alpha = a_0 + a_1 x/y + \ldots + a_n (x/y)^n =
b_0 + b_1 y/x + \ldots + b_m(y/x)^m
$$
for some $n, m \geq 0$ and $a_i, b_j \in R$.
Pick some $N > \max(n, m)$.
Consider the finite $R$-submodule $M$ of $R_{xy}$ generated by the elements
$$
(x/y)^N, (x/y)^{N - 1}, \ldots, x/y, 1, y/x, \ldots, (y/x)^{N - 1}, (y/x)^N
$$
We claim that $\alpha M \subset M$. Namely, it is clear that
$(x/y)^i (b_0 + b_1 y/x + \ldots + b_m(y/x)^m) \in M$ for
$0 \leq i \leq N$ and that
$(y/x)^i (a_0 + a_1 x/y + \ldots + a_n(x/y)^n) \in M$ for
$0 \leq i \leq N$. Hence $\alpha$ is integral over $R$ by
Lemma \ref{lemma-characterize-integral-element}. Note that
$\alpha \in R_x$, so if $R$ is integrally closed in $R_x$
then $\alpha \in R$ as desired.
\end{proof}